\chapter{Implementación}

En este capítulo veremos todo el proceso aplicado para implementar los diseños mostrados anteriormente. Veremos qué ha sido necesario hacer para poder construir todos los componentes que forman el sistema: el servidor back-end que hace de API, la \ac{BD}, el \ac{OS} y la interfaz de usuario.

\section{Servidor back-end}


\section{Base de datos}

Como ya hemos explicado, la \ac{BD} ha sido implementada con el motor PostgreSQL. Toda 

La \ac{BD} se ha implementado mediante los siguientes archivos .sql que dividen la creación de la \ac{BD} en cuatro pasos:

\begin{enumerate}
    \item \textbf{V1\_\_init\_schema}: crea todas las tablas necesarias para representar de forma correcta el modelo de datos.
    \item \textbf{V2\_\_add\_soft\_delete\_triggers}: crea los disparadores necesarios para poder aplicar un sistema de borrado lógico sobre los recursos cuya eliminación física acarrearía estados no deseados dentro del sistema.
    \item \textbf{V3\_\_create\_leaderboard\_view}: crea la vista encargada de calcular las puntuaciones de los equipos para cada liga, grupo y jornada.
    \item \textbf{V99\_\_insert\_initial\_data}: inserta los datos necesarios para que el sistema pueda funcionar de forma correcta: los deportes disponibles, roles de usuario y una configuración y sistema de puntuación por defecto para ligas de voleibol.
\end{enumerate}

% TODO: talk about leaderboard view table

Para la generación de los \ac{DTO} en todos los servicios se han utilizado herramientas específicas cuya función principal es la generación de clases a partir de un contrato API. En el back-end, se ha utilizado la librería \emph{org.openapi.generator}; y en el front-end, el paquete \emph{swagger-typescript-api}.

\section{Almacenamiento de objetos}


\section{Interfaz multiplataforma}

% Links to libraries:
% back: https://plugins.gradle.org/plugin/org.openapi.generator
% front: https://www.npmjs.com/package/swagger-typescript-api?activeTab=dependencies